\begin{table*}[!b]
\centering
\small
\setlength{\tabcolsep}{4.5pt}
\resizebox{\textwidth}{!}{%
\begin{tabular}{lrrrrrrr}
\toprule
Variant & Utility & $\Delta$ Utility & Action Acc. & $\Delta$ Acc. & Over-Answer & Answer-Supported Acc. & Defective-Premise Acc. \\
\midrule
baseline & -0.2188 & 0 & 0.775 & 0 & 0.05 & 0.7083 & 0.875 \\
decision\_first & \textbf{-0.1375} & +0.0813 & \textbf{0.85} & +0.075 & \textbf{0} & 0.75 & \textbf{1} \\
critique\_first & -0.2812 & -0.0624 & 0.4 & -0.375 & 0.025 & 0.0416 & 0.9375 \\
decision\_first\_guarded & -0.1812 & +0.0376 & 0.625 & -0.15 & \textbf{0} & 0.5 & 0.8125 \\
decision\_first\_balanced & -0.1313 & +0.0875 & 0.375 & -0.4 & \textbf{0} & 0 & 0.9375 \\
\bottomrule
\end{tabular}
}
\caption{Prompt-intervention pilot on the Qwen2.5-1.5B quick+stale split. Deltas are relative to the baseline prompt. Answer-supported accuracy pools answerable controls and conflicting-evidence items; defective-premise accuracy pools false- and stale-premise items. Utility is highlighted only among variants that do not reduce overall or answer-supported action accuracy versus baseline.}
\label{tab:qwen25-15b-intervention-pilot}
\end{table*}


This intervention is deliberately lightweight: it modifies instruction-level decision ordering rather than adding a new training stage or external verifier, following prompt-level calibration ideas from prior reasoning and self-consistency work\citep{wei-etal-2022-chain-of-thought,wang-etal-2023-self-consistency,madaan-etal-2023-self-refine}.

The intervention results suggest that action-selection failures are not fixed simply by adding more critique language. On Qwen2.5-1.5B quick+stale, \texttt{decision\_first} improves utility from \QwenQuickStaleInterventionBaselineUtility{} to \QwenQuickStaleInterventionDecisionFirstUtility{}, action accuracy from \QwenQuickStaleInterventionBaselineActionAcc{} to \QwenQuickStaleInterventionDecisionFirstActionAcc{}, and over-answer rate from \QwenQuickStaleInterventionBaselineOverAnswer{} to \QwenQuickStaleInterventionDecisionFirstOverAnswer{}. Crucially, this gain does not come from sacrificing answer-supported behavior: answer-supported accuracy changes from \QwenQuickStaleInterventionBaselineAnswerSupportedAcc{} to \QwenQuickStaleInterventionDecisionFirstAnswerSupportedAcc{}, while defective-premise accuracy rises from \QwenQuickStaleInterventionBaselineDefectivePremiseAcc{} to \QwenQuickStaleInterventionDecisionFirstDefectivePremiseAcc{}. In contrast, \texttt{critique\_first}, \texttt{decision\_first\_guarded}, and \texttt{decision\_first\_balanced} either over-challenge answerable items or collapse answer-supported accuracy. This supports a narrow method claim: making the action decision explicit can help, but longer or more cautious prompts are not automatically better.
